\section{Dataset Analysis}

\paragraph{Token count} At 14.7B tokens, OpenWebMath is just below the size of Minerva's Math Web Pages (17.5B tokens) \citet{lewkowycz2022solving} and significantly larger than the web part of any other dataset. OpenWebMath has around the same number of LLaMA tokens as ProofPile (14.2B) \citep{azerbayev2023proofnet}, but we note that there is very little overlap between between the two datasets. As a result, OpenWebMath brings a large number of new mathematical tokens that were previously unavailable to the open-source community. Due to differences in data curation strategies, it is hard to compare these datasets other than by training models on them. Since not much is known about how to properly filter a dataset, we opted to keep as much relevant content as possible. However, future work could explore filtering OpenWebMath more aggressively to further improve its quality.

\paragraph{Data Composition} We measured the distribution of domains in OpenWebMath both by document and by character count. \autoref{table:common-domains} and \autoref{table:top-domains-by-characters} show the top twenty most common domains by document and character count respectively. The most common sources of data tend to be discussion forums, blog posts, and scientific papers. We find that the distribution of characters in the dataset is distributed over 131,206 domains, with 46\% of the characters appearing in the top 100 domains.

\begin{table}[b!]
    \centering
    \begin{minipage}{0.48\textwidth}
        \centering
        \resizebox{\textwidth}{!}{
        %\parbox{9cm}{
        \begin{tabular}{lrr}
\multicolumn{1}{c}{\bf Domain}  &\multicolumn{1}{c}{\bf \# Documents} & \multicolumn{1}{c}{\bf \% Documents}\\ 
        \midrule
        stackexchange.com & 1,136,407 & 17.99\% \\ 
        physicsforums.com & 300,044 & 4.75\% \\ 
        mathhelpforum.com & 170,721 & 2.70\% \\ 
        socratic.org & 133,983 & 2.12\% \\ 
        mathoverflow.net & 120,755 & 1.91\% \\ 
        gradesaver.com & 96,100 & 1.52\% \\ 
        zbmath.org & 91,939 & 1.46\% \\ 
        wordpress.com & 87,876 & 1.39\% \\ 
        github.io & 81,125 & 1.28\% \\ 
        brilliant.org & 68,573 & 1.09\% \\ 
        gamedev.net & 50,560 & 0.80\% \\ 
        openstudy.com & 49,041 & 0.78\% \\ 
        gmatclub.com & 48,812 & 0.77\% \\ 
        blogspot.com & 48,036 & 0.76\% \\ 
        wikipedia.org & 46,606 & 0.74\% \\ 
        ac.uk & 41,342 & 0.65\% \\ 
        nature.com & 37,403 & 0.59\% \\ 
        aimsciences.org & 36,368 & 0.58\% \\ 
        libretexts.org & 32,216 & 0.51\% \\ 
        readthedocs.io & 31,455 & 0.50\% \\
        \end{tabular}
        }
        %}
        \caption{Most Common Domains by Document Count.}
        \label{table:common-domains}
    \end{minipage}\hfill
    \begin{minipage}{0.48\textwidth}
        \centering
        \resizebox{\textwidth}{!}{
        %\parbox{9cm}{
        \begin{tabular}{lrr}
        \multicolumn{1}{c}{\bf Domain}  &\multicolumn{1}{c}{\bf \# Characters} & \multicolumn{1}{c}{\bf \% Characters}\\
                \midrule
                stackexchange.com & 4,655,132,784 & 9.55\% \\ 
                nature.com & 1,529,935,838 & 3.14\% \\ 
                wordpress.com & 1,294,166,938 & 2.66\% \\ 
                physicsforums.com & 1,160,137,919 & 2.38\% \\ 
                github.io & 725,689,722 & 1.49\% \\ 
                zbmath.org & 620,019,503 & 1.27\% \\ 
                wikipedia.org & 618,024,754 & 1.27\% \\ 
                groundai.com & 545,214,990 & 1.12\% \\ 
                blogspot.com & 520,392,333 & 1.07\% \\ 
                mathoverflow.net & 499,102,560 & 1.02\% \\ 
                gmatclub.com & 442,611,169 & 0.91\% \\ 
                gamedev.net & 426,478,461 & 0.88\% \\ 
                ac.uk & 402,111,665 & 0.83\% \\ 
                aimsciences.org & 344,716,386 & 0.71\% \\ 
                mathhelpforum.com & 319,215,756 & 0.65\% \\ 
                deepai.org & 313,512,520 & 0.64\% \\ 
                libretexts.org & 282,014,149 & 0.58\% \\ 
                readthedocs.io & 269,816,413 & 0.55\% \\ 
                tib.eu & 199,714,017 & 0.41\% \\ 
                mit.edu & 198,487,362 & 0.41\% \\
        \end{tabular}
        }
        %}
        \caption{Most Common Domains by Character Count.}
        \label{table:top-domains-by-characters}
    \end{minipage}
\end{table}

In order to get a sense of the types of documents found in the dataset, we analyzed 100,000 randomly sampled documents. First, we created embeddings of this data using \texttt{all-MiniLM-L12-v2} \citep{wang2020minilm} in SentenceTransformers \citep{reimers-2019-sentence-bert}. Then, we clustered these embeddings using $k$-Means with $k=128$. Finally, we took the five closest documents to each cluster center and asked \texttt{gpt-3.5-turbo} (\href{https://platform.openai.com/docs/api-reference}{https://platform.openai.com/docs/api-reference}) to classify each cluster as Math, Physics, Statistics, Chemistry, Economics, Computer Science, or Other. We then aggregated these statistics, using the size of each cluster to get an estimate of the final number of documents in each category. We note several potential issues with this methodology, including inaccuracies stemming from using an LLM for classification, and the potential that not every document within a cluster belongs to the predicted category. \autoref{fig:doublefig} shows the results of this analysis. The majority of the documents in the dataset are directly related to mathematics, while the rest are spread out throughout physics, computer science, statistics, chemistry, and economics, with 12\% of documents not falling neatly into any of these categories.

We also used GPT to analyze the types of websites found in OpenWebMath. To do this, we took a sample of 200 documents and asked \texttt{gpt-3.5-turbo} to classify each as a Forum, Paper, Blog, Reference, Educational, Reference, or other. We also gave the document URL as a feature, since we found GPT is often able to judge the topic from the URL alone. We validated our analysis by asking GPT to do this classification on the top 100 domain names and got similar results. \autoref{fig:doublefig} shows the results. The highest proportion of documents are forum pages, where users ask and answer questions related to mathematical subjects. There is also a large proportion of educational and reference content.

\paragraph{Downstream Performance}
\label{sec:models}
We ran experiments to find out how our dataset compares to other language modeling datasets. We compare models trained on OpenWebMath for a single epoch (14.7B tokens) with models trained for the same number of tokens on The Pile \citep{gao2020pile}, a general langauge modeling dataset, and ProofPile \citep{azerbayev2023proofnet}, a dataset of both formal and informal mathematics. We also train a 50/50 mixture of ProofPile and OpenWebMath to evaluate the performance of OpenWebMath when included in a mixture of other datasets, as would be common in practice.

We train randomly initialized models with the same architecture as Pythia 1.4B \citep{pythia}. We use a batch size of 1M tokens and the same hyperparameters as Pythia otherwise. These models are evaluated on a collection of mathematics benchmarks which show signal on models of this size. This includes the subset of level-1 algebra questions from MATH, LILA-multiarith to test coding ability, and GSM8k and MATH perplexities, which scale more smoothly than accuracies. We also compare to Pythia 1.4B \citep{pythia}, which was trained on 300B tokens of The Pile \citep{gao2020pile} with the same architecture.

\autoref{table:math-perplexity-comparison} shows the results for our perplexity evaluations. There is a clear performance lead for models trained with OpenWebMath and the mixture seems to perform best. Despite Pythia being trained on over 20x the number of tokens, the performance of our models on the perplexity benchmarks far exceeds its performance, showing the potential of domain-specific models for mathematics. Similarly, \autoref{table:math-accuracy-comparison} shows the performance of the models on MATH-Algebra-Easy and LILA-multiarith \citep{lila}. OpenWebMath models outperform models that were not trained on it by a significant margin.